% Generado por sistema/backend ReporteTex.java -- no editar a mano.
% Datos: Base de datos: 148 importado (titulados) / Base de datos: 53 importado (empleadores)
\subsubsection*{C12. Edad $\times$ situación laboral}

\textit{Paso 1. Describir el problema.} Analizar el patrón demográfico: establecer si la situación laboral del titulado depende de su edad.

\begin{itemize}[nosep]
  \item Variable independiente $X$: Edad.
  \item Variable dependiente $Y$: Situación laboral.
  \item Población: todos los titulados ($N = 148$). Criterio ARCU-SUR: 3.2.3. Nivel de significancia $\alpha = 0.05$.
  \item $H_0$: Las variables «Edad» (X) y «Situación laboral» (Y) son independientes.
  \item $H_1$: Existe dependencia significativa entre las variables «Edad» (X) y «Situación laboral» (Y).
\end{itemize}

\textit{Datos.} Primeros 10 de los 148 titulados que respondieron ambas preguntas (sin datos personales).

\begin{mitabla}{C12: datos de los primeros titulados}
\begin{tabular}{lll}
\toprule
Titulado & $X$: Edad & $Y$: Situación laboral \\
\midrule
T1 & 27–34 & Emprendimiento propio \\
T2 & 27–34 & Trabaja en una organización \\
T3 & 35 o más & Trabaja en una organización \\
T4 & 26 o menos & Trabaja en una organización \\
T5 & 35 o más & Emprendimiento propio \\
T6 & 35 o más & Trabaja en una organización \\
T7 & 27–34 & Trabaja en una organización \\
T8 & 35 o más & Trabaja en una organización \\
T9 & 27–34 & No trabaja \\
T10 & 27–34 & Trabaja en una organización \\
\bottomrule
\end{tabular}
\end{mitabla}

\textit{Paso 2. Tabla de contingencia} (frecuencia observada $FO_{ij}$, $N = 148$). Cada celda cuenta a los titulados que respondieron a la vez la fila ($X$) y la columna ($Y$).

\begin{mitabla}{C12: tabla de contingencia, frecuencia observada $FO_{ij}$}
\begin{tabular}{L{3.8cm}C{3.00cm}C{3.00cm}C{3.00cm}|c}
\toprule
$X$: Edad \textbackslash{} $Y$: Situación laboral & Trabaja en una organización & Emprendimiento propio & No trabaja & Total \\
\midrule
26 o menos & 21 & 2 & 16 & 39 \\
27–34 & 35 & 6 & 4 & 45 \\
35 o más & 52 & 7 & 5 & 64 \\
\midrule
Total & 108 & 15 & 25 & $N = 148$ \\
\bottomrule
\end{tabular}
\end{mitabla}

\textit{Paso 3. Frecuencia esperada.} Es lo que habría en cada celda si $X$ e $Y$ fueran independientes. Por ejemplo, para la celda 26 o menos y Trabaja en una organización:
\[
FE_{11} = \frac{(39)(108)}{148} = 28.46
\]

\begin{mitabla}{C12: frecuencia esperada $FE_{ij}$}
\begin{tabular}{L{3.8cm}C{3.00cm}C{3.00cm}C{3.00cm}}
\toprule
$X$: Edad \textbackslash{} $Y$: Situación laboral & Trabaja en una organización & Emprendimiento propio & No trabaja \\
\midrule
26 o menos & $39\cdot108/148 = 28.46$ & $39\cdot15/148 = 3.95$ & $39\cdot25/148 = 6.59$ \\
27–34 & $45\cdot108/148 = 32.84$ & $45\cdot15/148 = 4.56$ & $45\cdot25/148 = 7.60$ \\
35 o más & $64\cdot108/148 = 46.70$ & $64\cdot15/148 = 6.49$ & $64\cdot25/148 = 10.81$ \\
\bottomrule
\end{tabular}
\end{mitabla}

\textit{Paso 4. Estadístico chi-cuadrado.} En cada celda se calcula $\frac{(FO_{ij}-FE_{ij})^2}{FE_{ij}}$ y luego se suman todas.

\begin{mitabla}{C12: cálculo de $X_0^2$ celda por celda}
\begin{tabular}{L{3.2cm}C{3.20cm}C{3.20cm}C{3.20cm}|c}
\toprule
$X$: Edad \textbackslash{} $Y$: Situación laboral & Trabaja en una organización & Emprendimiento propio & No trabaja & Suma fila \\
\midrule
26 o menos & $\dfrac{(21-28.46)^2}{28.46}$\newline $= 1.9552$ & $\dfrac{(2-3.95)^2}{3.95}$\newline $= 0.9647$ & $\dfrac{(16-6.59)^2}{6.59}$\newline $= 13.4473$ & $16.3672$ \\[6pt]
27–34 & $\dfrac{(35-32.84)^2}{32.84}$\newline $= 0.1424$ & $\dfrac{(6-4.56)^2}{4.56}$\newline $= 0.4541$ & $\dfrac{(4-7.60)^2}{7.60}$\newline $= 1.7062$ & $2.3027$ \\[6pt]
35 o más & $\dfrac{(52-46.70)^2}{46.70}$\newline $= 0.6009$ & $\dfrac{(7-6.49)^2}{6.49}$\newline $= 0.0407$ & $\dfrac{(5-10.81)^2}{10.81}$\newline $= 3.1233$ & $3.7648$ \\[6pt]
\bottomrule
\end{tabular}
\end{mitabla}

\[
X_0^2 = 16.3672 + 2.3027 + 3.7648 = 22.4347
\]

\textit{Paso 5. Estadístico teórico.} Con $f = 3$ filas y $c = 3$ columnas: $gl = (3-1)(3-1) = 4$. De la tabla chi-cuadrada, con $\alpha = 0.05$: $X^2_{4,\,0.05} = 9.488$.

\textit{Paso 6. Relación entre estadísticos y conclusión.}

\begin{center}
\begin{tabular}{l}
$22.4347 \ge 9.488 \rightarrow$ falso\\
Si ($X_0^2 < X^2_{(f-1)(c-1),\,\alpha}$)\\
\quad ``Independencia''\\
Sino\\
\quad \fbox{NO ``Independencia''}\\
FinSi
\end{tabular}
\end{center}

Se rechaza $H_0$ y se acepta $H_1$: «Situación laboral» depende de «Edad». Porcentaje de cada fila en la columna ``Trabaja en una organización'': 26 o menos $21/39 = 53.8\%$; 27–34 $35/45 = 77.8\%$; 35 o más $52/64 = 81.3\%$.

\textit{Observación:} 2 de 9 celdas tienen $FE < 5$, por lo que la prueba pierde confiabilidad y la conclusión debe tomarse con cautela.

\textit{Recomendación.} La situación laboral cambia con la edad: reforzar la vinculación con empleadores y la bolsa de trabajo para el grupo de edad con más titulados sin trabajo.

